\section{Calculated Sum-Frequency Spectra}

One of the aims of this simulation study is to provide complementary information to the SFG experimental studies of these systems.\cite{McFearin2009} The varied set of anions and their affect on the \ctc-\wat interface is linked from theory to empirical data through the connection of computed SFG spectra. The computed spectra for the SSP polarization are presented in figure \ref{fig:sfg-spectra} along with the experimental spectra (inserts) from our previous work with these same salt solutions interfaced with \ctc.\cite{McFearin2009} Each of the spectra show both a salt system response (colored traces) overlayed on the reference \ctcwat system spectrum (black or dashed-black traces). On first look, the overall computed intensities and lineshapes follow remarkably similar trends as in the experimental systems. All the spectra have a strong feature near 3660\cm coinciding with the ``free-OH'' vibrations as defined previously, and corresponding to one of the uncoupled OH stretch modes from water molecules that ``straddle'' the interface (figure \ref{fig:water-angles} a, b, and c).\cite{McFearin2009} The broad spectral region from 3200-3500\cm is attributed to the more highly-coordinated OH-oscillators that are solvated at the surface, or just beneath the surface with stronger hydrogen-bonding. Each of the spectra computed for the salt-solutions differ markedly from each other and from the neat \ctcwat spectrum. The monovalent ions (\cl and \nit) in solution produce a measurable decrease in intensity of the lower-frequencies of the spectra with very little change to the free-OH mode. Like the experimental counterparts, the decrease is greatest around the 3200-3400\cm region but shows little change from the neat-\ctcwat system above 3500\cm. Like in the SFG experiment, the presence of the \sul anion causes an opposite effect by significantly enhancing the spectrum below 3600\cm.

\begin{figure}[h!]
\begin{center}
	\includegraphics[scale=1.0]{images/sfg-spectra.png}
	\caption{Vibrational SFG spectra of the water-OH stretching region for each interfacial aqueous-salt-\ctc system. The reference \ctcwat interface spectrum (black-dashed line) is provided on each simulated SFG spectrum for reference. Inserts at the top left of each spectrum are reproductions of the experimental spectra from McFearin et al\cite{McFearin2009} including the reference (black) spectrum from that study. SSP polarization is used for all the spectra.}
	\label{fig:sfg-spectra}
\end{center}
\end{figure}

The reference \ctc-\wat spectrum reproduces well the lineshape from experiment, but lacks the definition of the two peaks found near 3250 and 3450\cm. These lower-frequency features have been attributed to the different H-bonding varieties of water that make up the more highly-coordinated, tetrahedral environments found deeper into the aqueous bulk. We have found difficulty in reproducing the lower-frequency peaks near 3200\cm and attribute the slight differences in the lineshape to similar problems reported previously.\cite{Walker2006b} However, the neat-\ctcwat lineshape is otherwise quite similar to that of the experiment. The salt solution signals spectra show an overall drop in signal from when \cl and \nit are added and an increase in intensity to \sul. This suggests that the methods are sound and justified for experimental comparison in this study.

%monovalent ion results
Our group's previous conclusions from experiment are that the presence of anions at the interface causes a ``field-screening'' that decreases the effect of the changing interfacial field further into the aqueous bulk. Water molecules found below the surface in the aqueous side, those giving rise to the lower-frequency (i.e. below 3500\cm) OH-stretch features, are thus less oriented by the field and their SFG contribution is decreased. Both the monovalent anions, \cl and \nit, show this effect in their SFG spectra. Comparison to the reference neat-\wat spectrum shows that the added presence of the surface-active anions decreases the lower-frequency intensities, as found in the experimental study. The \cl affects a notably smaller decrease in the spectral intensities than the \nit system. This expected result is most likely due to the higher preference for the surface of the larger, and more polarizable nitrate in the presence of \ctc. The \nit ion is extremely surface active, as seen in the density profile, and should thus cause the greatest ``field-screening'' to waters found deeper in the bulk.

% divalent/sulfate
The divalent and larger \sul anion accumulates deeper into the aqueous bulk and exhibits the lowest surface affinity of the ions studied. This is likely due to the higher charge of the anion that leads to greater solvation. The sulfate provides little or no screening of the interfacial field from the top-most water layer, and only affects the deeper waters. Lower frequency features of the water OH-spectrum spectrum are notably higher in intensity and both the 3250 and 3450\cm features are present, and dominate as compared to the reference water system, in both the simulated and experimental spectra.

As concluded in the previous experimental work, the monovalent anions appear to screen the deeper water molecules from the field produced by the phase change at the aqueous-organic interface. This is supported by the confirmation through our MD simulations here that monovalent ions show a strong surface affinity and interact with surface waters. The large but more highly charged divalent \sul anion experiences stronger solvation and is thus found deeper in the aqueous phase. Deeper anions do not participate as interfacial field screening agents to the same extent as their monovalent counterparts, but act to strongly orient the waters near the interface, perhaps through the strong field established by the ion double-layering. The distance between the counter-ion density peaks (table \ref{table:double-layer}) follows the inverse of the trend of SFG signal enhancement. As the ionic double-layer size increases, the SFG signal decreases. Similarly, the smallest double-layer size, that of the \sul system, produces the greatest signal enhancement across the lower frequencies of the water OH-stretching SFG spectrum. Although the water density profiles do not change markedly between the different systems, the orientational profiles do show large variation from the neat-\ctcwat system, and some slight variation between the salt-solutions. The two factors that may alter the SFG intensity of an experiment are changes in the number of contributing water bonded species, and a change in orientation of various water bonded species. From the water orientation profiles (figure \ref{fig:2dhisto}) it is clear that the presence of anions at the \ctcwat interface causes a strong orientation change from the reference \ctcwat system. Thus we find that there appears to be a strong coupling between the presence and size of an ionic double-layer, subsequent reorientation of surface water molecules, and the resulting SFG signal change.
